\cleardoublepage
\chapter*{Summary}
\addcontentsline{toc}{chapter}{Summary}

The development of multicore processors creates opportunities to accelerate computations through parallel processing. 
However, simply increasing the number of processing units does not guarantee a proportional increase in performance, 
as the results are influenced, among other factors, by the method of data partitioning, the organization of computations, 
and the overhead associated with process creation and synchronization.

The aim of this thesis was to analyze the performance and efficiency of selected parallel sorting algorithms on a multicore processor. 
Sequential and parallel implementations of Quick Sort, Merge Sort, Bucket Sort, and
Sample Sort were developed in Python using the \texttt{multiprocessing} module and shared memory.

The experiments were conducted on datasets ranging from 1,000 to 1,000,000 elements, 
including integer and floating-point data, random data, data with duplicates, and partially sorted data. 
Execution time, CPU and RAM usage, speedup, and efficiency were analyzed.

For a random dataset of 1,000,000 integers, the shortest execution time was obtained for Bucket Sort, 
at approximately 1.20~s with 16 processing units.
Quick Sort achieved its best result of approximately 1.80~s with 2 units, Merge Sort approximately 2.15~s with 4 units, 
and Sample Sort approximately 2.37~s with 8 units. 
The highest speedup was achieved by Sample Sort, reaching approximately 2.60, while for Bucket Sort, Merge Sort, 
and Quick Sort it was approximately 1.80, 1.59, and 1.21, respectively.

The results showed that increasing the number of processing units did not always lead to shorter execution times, 
while the achieved efficiency depended on the algorithm used, the dataset size, and the characteristics
of the input data. The greatest benefits from parallel processing were obtained for larger datasets. 
The conducted experiments confirmed that effective use of parallelism requires its level to be adapted to the
characteristics of the algorithm and the processed data.